\section{Validation and systematic checks}
\label{sec:validation}

The validation program asks whether the GPR model can interpolate the smooth prompt
background without absorbing a narrow signal, whether the fitted signal uncertainty is
calibrated, and whether the shared-$\eps^2$ combination behaves as expected. The
complete internal note contains many per-mass and per-sample diagnostic plots. The PRD
manuscript should keep only the main validation evidence in the body and leave the
full diagnostic grid to appendices or auxiliary material.

The first validation choice is the training exclusion around each tested mass.
Training the Gaussian process immediately outside a narrow $1.64\,\sigma_m$ window can
allow the near tails of an injected Gaussian signal to enter the sideband training set,
where they can be partly absorbed into the smooth background. Refmatched closure toys
therefore compare several training exclusions. The selected production geometry uses a
$2.25\,\sigma_m$ training exclusion matched to a $2.25\,\sigma_m$ blind and extraction
window. This choice reduces training-region signal leakage while retaining the signal
tails in the likelihood.

\begin{figure*}[t]
\centering
\begin{minipage}{0.47\textwidth}
\centering
\ifigure{\linewidth}{blind225_schematic.png}\\
(a) Selected $2.25\,\sigma_m$ blind/extraction window.
\end{minipage}
\hfill
\begin{minipage}{0.47\textwidth}
\centering
\ifigure{\linewidth}{blind_geometry_reach.png}\\
(b) Refmatched blind-geometry reach comparison.
\end{minipage}
\caption{Training-exclusion and blind-window validation. The selected
$2.25\,\sigma_m$ geometry excludes nearby signal tails from the GPR training set while
including those tails in the signal extraction likelihood. The reach comparison shows
that the widened blind window preserves closure behavior while improving the expected
signal-yield reach relative to discarding the $1.64$--$2.25\,\sigma_m$ tails as a
guard band.}
\label{fig:blind-validation}
\end{figure*}

The second validation layer is closure on pseudoexperiments. Smooth analytic
functional-form toys provide a background model independent of the Gaussian process
used in the extraction. Full-refit toys then rerun the sideband-trained GPR procedure
on pseudo-data with known injected signal strengths. The principal metrics are the
signal-yield pull mean, pull width, recovered-over-injected yield, and residual
Gaussian-equivalent significance. Pull means near zero test bias; pull widths near one
test the fitted uncertainty. Pull widths below one, when the toys are faithful, imply
conservative uncertainties rather than false exclusions.

\begin{figure*}[t]
\centering
\ifigure{0.88\textwidth}{toy_upper_limit_proxy.png}
\caption{Three-sample toy upper-limit proxy summary for the 2015, 2016 10\%, and 2021
1\% validation samples. These background-only pseudoexperiment studies check the
expected signal-yield and $\eps^2$ scales before interpreting the observed
shared-coupling result.}
\label{fig:toy-upper-proxy}
\end{figure*}

The third layer checks sensitivity to systematic modeling choices. These include the
mass-resolution parameterization, the radiative-fraction normalization, the effective
length-scale bounds, the 2021 target-constrained smearing scale, and the confidence
level convention used for external comparisons. The present draft includes the
corresponding figures in the appendix. For the final journal article, systematic
variations should be promoted to the body only if they enter the quoted result or
materially change the external contour.

Finally, discovery calibration is kept distinct from upper-limit diagnostics. The
limit-tail plots compare observed upper limits with background-only limit ensembles and
answer whether an observed limit is unusually strong or weak. They are not discovery
$p$-values. A discovery statement must instead use the one-sided local $p_0$ statistic
and a scan-level look-elsewhere calibration. The current note uses effective-trials
overlays as review guides; a final PRD discovery statement should use a frozen toy
calibration of $q_{0,\max}$ over the final scan range.
